%====================================================
\section{Annales officielles et sujets de brevet cliquables}

\begin{objectif}
Cette section regroupe des liens cliquables vers des sujets de brevet de mathématiques des années précédentes. 
Le but est de travailler sur de vrais sujets, puis de compléter avec des sujets inventés plus difficiles.
\end{objectif}

\begin{retenir}
\textbf{Banques d'annales à utiliser régulièrement :}
\begin{itemize}
  \item \href{https://www.apmep.fr/Annales-du-Brevet-des-colleges}{APMEP -- Annales du brevet des collèges}.
  \item \href{https://www.apmep.fr/Brevet-2025}{APMEP -- Brevet 2025 et sujets zéro 2026}.
  \item \href{https://www.apmep.fr/Brevet-2024}{APMEP -- Brevet 2024}.
  \item \href{https://www.apmep.fr/Brevet-2023}{APMEP -- Brevet 2023}.
  \item \href{https://www.apmep.fr/Brevet-2022}{APMEP -- Brevet 2022}.
  \item \href{https://www.apmep.fr/Brevet-2021}{APMEP -- Brevet 2021}.
  \item \href{https://www.lumni.fr/article/maths-sujets-et-corriges-serie-generale-brevet-2025}{Lumni -- Maths brevet 2025 série générale, sujet et corrigé}.
  \item \href{https://www.letudiant.fr/college/brevet/sujets-et-corriges-du-brevet-de-mathematiques-2025-dnb.html}{L'Étudiant -- Sujet et corrigé du brevet de mathématiques 2025}.
  \item \href{https://groupe-reussite.fr/ressources/annales-brevet-maths/}{Groupe Réussite -- Annales de brevet de mathématiques}.
\end{itemize}
\end{retenir}

\begin{methode}
\textbf{Méthode pour travailler une annale}
\begin{enumerate}
  \item Faire le sujet en temps limité, sans regarder la correction.
  \item Noter le temps passé sur chaque exercice.
  \item Corriger en distinguant les erreurs de cours, de calcul, de méthode et de rédaction.
  \item Refaire les exercices ratés deux jours plus tard.
  \item Recopier dans une fiche les formules ou méthodes oubliées.
\end{enumerate}
\end{methode}

%====================================================
\section{Programme de révision jusqu'au brevet}

\begin{objectif}
Le planning suivant est conçu pour aller jusqu'à l'épreuve de mathématiques du brevet. Il mélange :
\begin{itemize}
  \item révision ciblée du cours ;
  \item exercices inventés plus difficiles ;
  \item sujets de brevet des années précédentes ;
  \item correction active ;
  \item révisions de physique-chimie pour l'épreuve de sciences.
\end{itemize}
\end{objectif}

\renewcommand{\arraystretch}{1.35}

\begin{longtable}{|p{2.3cm}|p{3.2cm}|p{4.7cm}|p{4.7cm}|p{0.7cm}|}
\hline
\textbf{Date} & \textbf{Objectif} & \textbf{Mathématiques} & \textbf{Physique-Chimie} & \textbf{Fait} \\
\hline
Ven. 29 mai & Diagnostic initial & Calculs courts : fractions, puissances, priorités. Puis sujet inventé \textbf{M-LONG-01, partie A}. & Conversions : km/h, m/s, L, mL, g, kg. Exercice PC-LONG-01, partie A. & $\square$ \\
\hline
Sam. 30 mai & Calcul numérique & Fractions, puissances de 10, notation scientifique, racines carrées. Puis \textbf{M-LONG-01, partie B}. & Mouvement, vitesse moyenne, distance, durée. PC-LONG-01, partie B. & $\square$ \\
\hline
Dim. 31 mai & Correction active & Refaire toutes les erreurs du diagnostic. Faire une fiche : calculs dangereux. & Refaire les conversions et formules ratées. & $\square$ \\
\hline
Lun. 1 juin & Calcul littéral & Développer, réduire, factoriser. Puis \textbf{M-LONG-01, partie C}. & Électricité : tension, intensité, résistance, loi d'Ohm. & $\square$ \\
\hline
Mar. 2 juin & Équations & Équations du premier degré, équations produit, mise en équation. Faire \textbf{M-LONG-01 en entier}. & Puissance et énergie électrique. & $\square$ \\
\hline
Mer. 3 juin & Fonctions & Images, antécédents, lecture graphique, fonctions affines. Commencer \textbf{M-LONG-02, partie A}. & Signaux : son, lumière, vitesse de propagation. & $\square$ \\
\hline
Jeu. 4 juin & Proportionnalité & Pourcentages, échelles, coefficients multiplicateurs, vitesse. Faire \textbf{M-LONG-02, partie B}. & Masse volumique, densité, identification d'un matériau. & $\square$ \\
\hline
Ven. 5 juin & Géométrie plane & Pythagore, réciproque de Pythagore, trigonométrie. Faire \textbf{M-LONG-02, partie C}. & Atomes, ions, molécules. & $\square$ \\
\hline
Sam. 6 juin & Thalès et homothéties & Thalès, réciproque, triangles semblables, agrandissement-réduction. Faire \textbf{M-LONG-02 en entier}. & Transformations chimiques et équations équilibrées. & $\square$ \\
\hline
Dim. 7 juin & Annale 2025 & Faire un sujet sur \href{https://www.apmep.fr/Brevet-2025}{APMEP -- Brevet 2025}. Durée : 2 h. & Relecture fiches sciences. & $\square$ \\
\hline
Lun. 8 juin & Correction annale & Corriger l'annale 2025. Refaire deux exercices ratés sans correction. & Refaire les questions d'unités ratées. & $\square$ \\
\hline
Mar. 9 juin & Statistiques & Moyenne, médiane, étendue, fréquences. Commencer \textbf{M-LONG-03, partie A}. & Énergie : joule, kilowattheure, coût. & $\square$ \\
\hline
Mer. 10 juin & Probabilités & Probabilités simples, événement contraire, tableaux à deux entrées. Faire \textbf{M-LONG-03, partie B}. & Loi d'Ohm et circuits simples. & $\square$ \\
\hline
Jeu. 11 juin & Algorithmique & Programmes Scratch/Python, boucles, variables, seuils. Faire \textbf{M-LONG-03, partie C}. & Signaux : écho, distance, vitesse du son. & $\square$ \\
\hline
Ven. 12 juin & Annale 2024 & Faire un sujet sur \href{https://www.apmep.fr/Brevet-2024}{APMEP -- Brevet 2024}. Durée : 2 h. & Mini-sujet sciences inventé. & $\square$ \\
\hline
Sam. 13 juin & Correction annale & Corriger l'annale 2024. Remplir le carnet d'erreurs. & Refaire les exercices sciences ratés. & $\square$ \\
\hline
Dim. 14 juin & Repos actif & 45 min : formules essentielles et erreurs classiques. & 30 min : unités et équations chimiques. & $\square$ \\
\hline
Lun. 15 juin & Sujet inventé difficile & Faire \textbf{M-LONG-03 en entier} en temps limité. & Faire \textbf{PC-LONG-01 en entier}. & $\square$ \\
\hline
Mar. 16 juin & Géométrie espace & Volumes, sphère, cylindre, cône, sections, agrandissement de volumes. Faire \textbf{M-LONG-04, partie A}. & Masse volumique et énergie. & $\square$ \\
\hline
Mer. 17 juin & Annale 2023 & Faire un sujet sur \href{https://www.apmep.fr/Brevet-2023}{APMEP -- Brevet 2023}. Durée : 2 h. & Relecture chimie. & $\square$ \\
\hline
Jeu. 18 juin & Correction annale & Correction active de l'annale 2023. Refaire les questions perdues. & Électricité + énergie. & $\square$ \\
\hline
Ven. 19 juin & Sujet transversal & Faire \textbf{M-LONG-04, parties B et C}. & Faire \textbf{PC-LONG-02, partie A}. & $\square$ \\
\hline
Sam. 20 juin & Brevet blanc maison & Faire \textbf{BB-MATHS-01}, sujet complet inventé difficile. Durée : 2 h. & Faire \textbf{BB-PC-01}. Durée : 1 h. & $\square$ \\
\hline
Dim. 21 juin & Correction brevet blanc & Corriger le brevet blanc. Refaire les exercices ratés. & Corriger le sujet sciences. & $\square$ \\
\hline
Lun. 22 juin & Fonctions et équations & Refaire tous les exercices de fonctions et équations ratés. & Revoir loi d'Ohm et énergie. & $\square$ \\
\hline
Mar. 23 juin & Problèmes complexes & Faire \textbf{M-LONG-05}, sujet transversal sur un voyage scolaire. & Faire \textbf{PC-LONG-02, partie B}. & $\square$ \\
\hline
Mer. 24 juin & Annale 2022 & Faire un sujet sur \href{https://www.apmep.fr/Brevet-2022}{APMEP -- Brevet 2022}. Durée : 2 h. & Mini-sujet sciences. & $\square$ \\
\hline
Jeu. 25 juin & Correction finale & Corriger l'annale 2022. Ne plus découvrir de nouveau chapitre. & Relecture complète sciences. & $\square$ \\
\hline
Ven. 26 juin & Épreuve de français & Maths léger : 30 min de formules seulement. & Sciences léger : unités et schémas. & $\square$ \\
\hline
Sam. 27 juin & Dernier sujet & Faire un dernier sujet représentatif sur \href{https://www.lumni.fr/article/maths-sujets-et-corriges-serie-generale-brevet-2025}{Lumni -- Brevet maths 2025}. & Revoir physique-chimie. & $\square$ \\
\hline
Dim. 28 juin & Repos intelligent & Refaire uniquement les erreurs déjà corrigées. Pas de nouveau sujet difficile. & Dernière fiche sciences. & $\square$ \\
\hline
Lun. 29 juin & Épreuve de sciences & Le soir : 45 min de formules maths, pas plus. & Épreuve de sciences. & $\square$ \\
\hline
Mar. 30 juin & Épreuve de maths & Brevet de mathématiques : commencer par le plus rentable, rédiger, relire. & Repos après l'épreuve. & $\square$ \\
\hline
\end{longtable}

\newpage

%====================================================
\section{Sujets inventés longs et difficiles de mathématiques}

\begin{attention}
Les sujets suivants sont inventés. Ils sont plus longs et parfois plus difficiles qu'un exercice classique de brevet. 
Ils servent à préparer les bons élèves, à travailler la prise d'initiative et à consolider toutes les méthodes du programme de troisième.
\end{attention}

%====================================================
\subsection{M-LONG-01 -- Calculs, programme de calcul et équations}

\begin{objectif}
Durée conseillée : 1 h.\\
Chapitres : fractions, puissances, calcul littéral, équations, équations produit.
\end{objectif}

\begin{exercice}[M-LONG-01 -- Partie A : calcul numérique]
On considère les nombres :
\[
A=7-3(2-5)^2+\frac{5}{6}\div\frac{10}{9},
\]
\[
B=\frac{2^5\times 2^{-3}}{4}+\frac{3}{5}-\frac{7}{10},
\]
\[
C=\left(\frac{3}{4}-\frac{5}{6}\right)^2.
\]

\begin{enumerate}
  \item Calculer $(2-5)^2$.
  \item Calculer $\dfrac{5}{6}\div\dfrac{10}{9}$.
  \item En déduire la valeur exacte de $A$.
  \item Simplifier $2^5\times2^{-3}$.
  \item Calculer $B$ sous forme d'une fraction irréductible.
  \item Calculer $C$ sous forme d'une fraction irréductible.
  \item Ranger $A$, $B$ et $C$ dans l'ordre croissant.
\end{enumerate}
\end{exercice}

\begin{exercice}[M-LONG-01 -- Partie B : programme de calcul]
On considère le programme de calcul suivant :
\begin{enumerate}
  \item Choisir un nombre.
  \item Ajouter $4$.
  \item Multiplier le résultat par le nombre choisi au départ.
  \item Soustraire le carré du nombre choisi.
  \item Ajouter $7$.
\end{enumerate}

\begin{enumerate}
  \item Appliquer ce programme au nombre $3$.
  \item Appliquer ce programme au nombre $-5$.
  \item On note $x$ le nombre choisi au départ. Montrer que le résultat final est $4x+7$.
  \item Quel nombre faut-il choisir pour obtenir $51$ ?
  \item Peut-on obtenir $-20$ ? Justifier.
\end{enumerate}
\end{exercice}

\begin{exercice}[M-LONG-01 -- Partie C : expression factorisée]
On considère :
\[
E=(x+4)^2-(x+4)(2x-1).
\]

\begin{enumerate}
  \item Développer et réduire $E$.
  \item Factoriser $E$ sans développer.
  \item Résoudre l'équation $E=0$.
  \item Résoudre l'équation $E=9$ en utilisant la forme développée.
  \item Vérifier une solution obtenue à la question précédente.
\end{enumerate}
\end{exercice}

\begin{correction}[M-LONG-01 -- Correction synthétique]
Pour la partie A :
\[
(2-5)^2=9.
\]
\[
\frac{5}{6}\div\frac{10}{9}=\frac{5}{6}\times\frac{9}{10}=\frac{45}{60}=\frac34.
\]
Donc :
\[
A=7-27+\frac34=-20+\frac34=-\frac{77}{4}.
\]
Pour $B$ :
\[
2^5\times2^{-3}=2^2=4,
\]
donc :
\[
B=\frac44+\frac35-\frac7{10}=1+\frac6{10}-\frac7{10}=\frac9{10}.
\]
Pour le programme :
\[
x\rightarrow x+4\rightarrow x(x+4)=x^2+4x,
\]
puis :
\[
x^2+4x-x^2+7=4x+7.
\]
Pour obtenir $51$ :
\[
4x+7=51,\quad 4x=44,\quad x=11.
\]
Pour $E$ :
\[
E=(x+4)\left[(x+4)-(2x-1)\right]=(x+4)(5-x).
\]
Donc :
\[
E=0 \Longleftrightarrow x=-4 \text{ ou } x=5.
\]
\end{correction}

\newpage

%====================================================
\subsection{M-LONG-02 -- Fonctions, tarifs et proportionnalité}

\begin{objectif}
Durée conseillée : 1 h 15.\\
Chapitres : fonctions affines, proportionnalité, pourcentages, inéquations, interprétation.
\end{objectif}

\begin{exercice}[M-LONG-02 -- Partie A : deux tarifs]
Une ville propose deux tarifs pour louer un vélo électrique :
\begin{itemize}
  \item Tarif A : $0{,}30$ euro par minute.
  \item Tarif B : abonnement de $5$ euros puis $0{,}12$ euro par minute.
\end{itemize}

On note $x$ le nombre de minutes de location.

\begin{enumerate}
  \item Écrire l'expression $A(x)$ du prix avec le tarif A.
  \item Écrire l'expression $B(x)$ du prix avec le tarif B.
  \item Calculer $A(20)$ et $B(20)$.
  \item Calculer $A(45)$ et $B(45)$.
  \item Résoudre l'inéquation $B(x)<A(x)$.
  \item À partir de quelle durée entière le tarif B devient-il plus avantageux ?
\end{enumerate}
\end{exercice}

\begin{exercice}[M-LONG-02 -- Partie B : lecture et interprétation]
On considère les fonctions :
\[
A(x)=0{,}30x
\quad \text{et} \quad
B(x)=5+0{,}12x.
\]

\begin{enumerate}
  \item Dire si $A$ est une fonction linéaire ou affine.
  \item Dire si $B$ est une fonction linéaire ou affine.
  \item Que représente le coefficient $0{,}30$ dans le contexte ?
  \item Que représente le nombre $5$ dans l'expression de $B(x)$ ?
  \item Résoudre $A(x)=B(x)$.
  \item Interpréter graphiquement le résultat.
\end{enumerate}
\end{exercice}

\begin{exercice}[M-LONG-02 -- Partie C : augmentation de prix]
L'année suivante, la ville augmente :
\begin{itemize}
  \item le tarif A de $10\%$ ;
  \item l'abonnement du tarif B de $20\%$ ;
  \item le prix par minute du tarif B de $5\%$.
\end{itemize}

\begin{enumerate}
  \item Donner le nouveau prix par minute du tarif A.
  \item Donner le nouveau prix de l'abonnement du tarif B.
  \item Donner le nouveau prix par minute du tarif B.
  \item Écrire les nouvelles fonctions $A_2(x)$ et $B_2(x)$.
  \item Pour $40$ minutes, comparer les deux nouveaux tarifs.
  \item Résoudre $B_2(x)<A_2(x)$.
\end{enumerate}
\end{exercice}

\begin{correction}[M-LONG-02 -- Correction synthétique]
\[
A(x)=0{,}30x,\qquad B(x)=5+0{,}12x.
\]
Pour $20$ minutes :
\[
A(20)=6,\qquad B(20)=7{,}4.
\]
Pour $45$ minutes :
\[
A(45)=13{,}5,\qquad B(45)=10{,}4.
\]
On résout :
\[
5+0{,}12x<0{,}30x,
\]
\[
5<0{,}18x,
\]
\[
x>\frac{5}{0{,}18}\approx 27{,}78.
\]
Le tarif B devient plus avantageux à partir de $28$ minutes.

Après augmentation :
\[
0{,}30\times1{,}10=0{,}33.
\]
\[
5\times1{,}20=6.
\]
\[
0{,}12\times1{,}05=0{,}126.
\]
Donc :
\[
A_2(x)=0{,}33x,\qquad B_2(x)=6+0{,}126x.
\]
\end{correction}

\newpage

%====================================================
\subsection{M-LONG-03 -- Statistiques, probabilités et algorithmique}

\begin{objectif}
Durée conseillée : 1 h 20.\\
Chapitres : moyenne, médiane, étendue, probabilités, simulation, algorithmique.
\end{objectif}

\begin{exercice}[M-LONG-03 -- Partie A : statistiques]
Une application de révision enregistre le nombre de questions réussies par 15 élèves lors d'un entraînement :
\[
12,\ 15,\ 9,\ 18,\ 20,\ 14,\ 16,\ 10,\ 11,\ 13,\ 19,\ 17,\ 15,\ 8,\ 14.
\]

\begin{enumerate}
  \item Calculer l'étendue.
  \item Calculer la moyenne.
  \item Ranger les valeurs dans l'ordre croissant.
  \item Déterminer la médiane.
  \item Combien d'élèves ont réussi au moins 15 questions ?
  \item Calculer la fréquence correspondante.
\end{enumerate}
\end{exercice}

\begin{exercice}[M-LONG-03 -- Partie B : probabilités]
L'application propose ensuite un quiz aléatoire. Une question est choisie parmi :
\begin{itemize}
  \item 5 questions de calcul ;
  \item 4 questions de géométrie ;
  \item 3 questions de probabilités ;
  \item 2 questions de physique.
\end{itemize}

\begin{enumerate}
  \item Combien y a-t-il de questions au total ?
  \item Quelle est la probabilité de tomber sur une question de géométrie ?
  \item Quelle est la probabilité de ne pas tomber sur une question de physique ?
  \item On tire une question, puis une deuxième avec remise. Quelle est la probabilité d'obtenir deux questions de calcul ?
  \item Quelle est la probabilité d'obtenir au moins une question de physique en deux tirages avec remise ?
\end{enumerate}
\end{exercice}

\begin{exercice}[M-LONG-03 -- Partie C : algorithme de seuil]
Un élève révise avec l'application. Il gagne $18$ points le premier jour. Chaque jour suivant, il gagne $4$ points de plus que la veille.

On considère l'algorithme suivant :

\begin{verbatim}
total = 0
jour = 0
points = 18

tant que total < 500:
    total = total + points
    jour = jour + 1
    points = points + 4

afficher(jour)
\end{verbatim}

\begin{enumerate}
  \item Que vaut \texttt{total} après le premier jour ?
  \item Que vaut \texttt{points} au début du deuxième jour ?
  \item Compléter le tableau suivant jusqu'au jour 5.

\begin{center}
\begin{tabular}{|c|c|c|}
\hline
Jour & Points gagnés ce jour & Total \\
\hline
1 & 18 & 18 \\
\hline
2 & & \\
\hline
3 & & \\
\hline
4 & & \\
\hline
5 & & \\
\hline
\end{tabular}
\end{center}

  \item Que représente la variable \texttt{jour} affichée à la fin ?
  \item Par essais successifs, déterminer le premier jour où le total dépasse 500 points.
\end{enumerate}
\end{exercice}

\begin{correction}[M-LONG-03 -- Correction synthétique]
L'étendue vaut :
\[
20-8=12.
\]
La somme des valeurs vaut :
\[
211.
\]
La moyenne est :
\[
\frac{211}{15}\approx14{,}1.
\]
Valeurs rangées :
\[
8,9,10,11,12,13,14,14,15,15,16,17,18,19,20.
\]
Il y a 15 valeurs, donc la médiane est la 8e valeur :
\[
14.
\]
Pour les probabilités, il y a :
\[
5+4+3+2=14
\]
questions.
\[
P(\text{géométrie})=\frac{4}{14}=\frac27.
\]
\[
P(\text{pas physique})=\frac{12}{14}=\frac67.
\]
Deux calculs avec remise :
\[
\frac{5}{14}\times\frac{5}{14}=\frac{25}{196}.
\]
Au moins une physique :
\[
1-P(\text{aucune physique})
=
1-\frac{12}{14}\times\frac{12}{14}
=
1-\frac{36}{49}
=
\frac{13}{49}.
\]
\end{correction}

\newpage

%====================================================
\subsection{M-LONG-04 -- Géométrie, volumes et optimisation}

\begin{objectif}
Durée conseillée : 1 h 30.\\
Chapitres : Pythagore, trigonométrie, volumes, agrandissement, pourcentages.
\end{objectif}

\begin{exercice}[M-LONG-04 -- Partie A : rampe d'accès]
Une rampe d'accès est modélisée par un triangle rectangle. La longueur horizontale est $4{,}80$ m et la hauteur est $0{,}90$ m.

\begin{enumerate}
  \item Faire un schéma.
  \item Calculer la longueur de la rampe.
  \item Calculer la valeur approchée de l'angle formé avec le sol.
  \item Une norme impose un angle inférieur à $12^\circ$. La rampe respecte-t-elle la norme ?
\end{enumerate}
\end{exercice}

\begin{exercice}[M-LONG-04 -- Partie B : réservoir cylindrique]
Un réservoir d'eau a la forme d'un cylindre de rayon $0{,}75$ m et de hauteur $1{,}80$ m.

\begin{enumerate}
  \item Calculer le volume du réservoir en $\mathrm{m^3}$.
  \item Convertir ce volume en litres, sachant que $1\ \mathrm{m^3}=1000\ \mathrm{L}$.
  \item Le réservoir est rempli à $80\%$. Quel volume d'eau contient-il ?
  \item Une famille consomme $150$ L par jour. Pendant combien de jours peut-elle utiliser ce réservoir ?
\end{enumerate}
\end{exercice}

\begin{exercice}[M-LONG-04 -- Partie C : agrandissement]
On fabrique une maquette du réservoir avec un coefficient de réduction de $0{,}2$.

\begin{enumerate}
  \item Par combien les longueurs sont-elles multipliées ?
  \item Par combien les aires sont-elles multipliées ?
  \item Par combien les volumes sont-ils multipliés ?
  \item Calculer le volume de la maquette.
  \item Expliquer pourquoi le volume n'est pas multiplié par $0{,}2$.
\end{enumerate}
\end{exercice}

\begin{correction}[M-LONG-04 -- Correction synthétique]
Pour la rampe :
\[
L^2=4{,}80^2+0{,}90^2=23{,}04+0{,}81=23{,}85.
\]
\[
L=\sqrt{23{,}85}\approx4{,}88\ \mathrm{m}.
\]
Pour l'angle $\alpha$ :
\[
\tan(\alpha)=\frac{0{,}90}{4{,}80}=0{,}1875.
\]
Donc :
\[
\alpha\approx10{,}6^\circ.
\]
La rampe respecte la norme.

Pour le cylindre :
\[
V=\pi r^2h=\pi\times0{,}75^2\times1{,}80.
\]
\[
V\approx3{,}18\ \mathrm{m^3}.
\]
Donc environ :
\[
3180\ \mathrm{L}.
\]
À $80\%$ :
\[
3180\times0{,}80=2544\ \mathrm{L}.
\]
Nombre de jours :
\[
\frac{2544}{150}\approx16{,}96.
\]
La famille peut tenir environ 16 jours complets.

Pour la maquette, les volumes sont multipliés par :
\[
0{,}2^3=0{,}008.
\]
\end{correction}

\newpage

%====================================================
\subsection{M-LONG-05 -- Problème transversal : voyage scolaire}

\begin{objectif}
Durée conseillée : 1 h 45.\\
Chapitres : proportionnalité, fonctions, statistiques, probabilités, vitesse, calcul littéral.
\end{objectif}

\begin{exercice}[M-LONG-05 -- Partie A : coût du voyage]
Un collège organise un voyage scolaire. Le coût total dépend du nombre $x$ d'élèves participants.

Le transport coûte $720$ euros au total. L'hébergement coûte $38$ euros par élève. Les visites coûtent $12$ euros par élève.

\begin{enumerate}
  \item Exprimer le coût total $C(x)$ du voyage.
  \item Calculer $C(40)$.
  \item Calculer le coût par élève si 40 élèves participent.
  \item Le collège souhaite que le coût par élève ne dépasse pas $70$ euros. Écrire une inéquation.
  \item Déterminer le nombre minimal d'élèves nécessaires.
\end{enumerate}
\end{exercice}

\begin{exercice}[M-LONG-05 -- Partie B : trajet]
Le bus parcourt $315$ km. Il roule à une vitesse moyenne de $90\ \mathrm{km/h}$ pendant les deux premières heures, puis à $75\ \mathrm{km/h}$ pour le reste du trajet.

\begin{enumerate}
  \item Quelle distance parcourt-il pendant les deux premières heures ?
  \item Quelle distance reste-t-il à parcourir ?
  \item Combien de temps dure la deuxième partie du trajet ?
  \item Calculer la durée totale du trajet.
  \item Calculer la vitesse moyenne sur tout le trajet.
\end{enumerate}
\end{exercice}

\begin{exercice}[M-LONG-05 -- Partie C : chambres]
L'auberge propose des chambres de 3 et de 4 élèves. Il y a 41 élèves.

\begin{enumerate}
  \item Peut-on utiliser uniquement des chambres de 4 ? Justifier.
  \item Trouver une répartition possible avec des chambres de 3 et de 4.
  \item On note $a$ le nombre de chambres de 3 et $b$ le nombre de chambres de 4. Écrire une équation vérifiée par $a$ et $b$.
  \item Trouver deux solutions possibles.
  \item Quelle solution utilise le moins de chambres ?
\end{enumerate}
\end{exercice}

\begin{exercice}[M-LONG-05 -- Partie D : probabilités]
Pendant le voyage, un professeur choisit au hasard un élève pour répondre à une question. Parmi les 41 élèves, 24 sont des filles. Parmi les élèves, 10 ont déjà visité le musée, dont 6 filles.

\begin{enumerate}
  \item Quelle est la probabilité de choisir une fille ?
  \item Quelle est la probabilité de choisir un élève ayant déjà visité le musée ?
  \item Quelle est la probabilité de choisir une fille ayant déjà visité le musée ?
  \item Quelle est la probabilité de choisir un garçon n'ayant jamais visité le musée ?
\end{enumerate}
\end{exercice}

\begin{correction}[M-LONG-05 -- Correction synthétique]
Coût :
\[
C(x)=720+38x+12x=720+50x.
\]
Pour 40 élèves :
\[
C(40)=720+50\times40=2720.
\]
Coût par élève :
\[
\frac{2720}{40}=68.
\]
Condition :
\[
\frac{720+50x}{x}\leq70.
\]
\[
720+50x\leq70x.
\]
\[
720\leq20x.
\]
\[
x\geq36.
\]
Il faut au moins 36 élèves.

Trajet :
\[
90\times2=180\ \mathrm{km}.
\]
Il reste :
\[
315-180=135\ \mathrm{km}.
\]
Durée de la deuxième partie :
\[
\frac{135}{75}=1{,}8\ \mathrm{h}.
\]
Durée totale :
\[
3{,}8\ \mathrm{h}.
\]
Vitesse moyenne :
\[
\frac{315}{3{,}8}\approx82{,}9\ \mathrm{km/h}.
\]
Chambres :
\[
3a+4b=41.
\]
Une solution est :
\[
a=3,\quad b=8,
\]
car :
\[
3\times3+4\times8=9+32=41.
\]
Probabilités :
\[
P(\text{fille})=\frac{24}{41}.
\]
\[
P(\text{déjà visité})=\frac{10}{41}.
\]
\[
P(\text{fille et déjà visité})=\frac{6}{41}.
\]
Garçons : $41-24=17$. Garçons ayant déjà visité : $10-6=4$. Garçons n'ayant jamais visité :
\[
17-4=13.
\]
Donc :
\[
P=\frac{13}{41}.
\]
\end{correction}

\newpage

%====================================================
\section{Sujets inventés longs de physique-chimie}

\subsection{PC-LONG-01 -- Vélo électrique, énergie et vitesse}

\begin{objectif}
Durée conseillée : 1 h.\\
Chapitres : vitesse, énergie, puissance, conversions, autonomie.
\end{objectif}

\begin{exercice}[PC-LONG-01 -- Partie A : vitesse]
Un élève utilise un vélo électrique pour parcourir $18$ km. Il met $54$ minutes.

\begin{enumerate}
  \item Convertir $54$ minutes en heures.
  \item Calculer la vitesse moyenne en $\mathrm{km/h}$.
  \item Convertir cette vitesse en $\mathrm{m/s}$.
  \item La vitesse maximale autorisée avec assistance est $25\ \mathrm{km/h}$. Le vélo respecte-t-il cette limite en moyenne ?
\end{enumerate}
\end{exercice}

\begin{exercice}[PC-LONG-01 -- Partie B : énergie consommée]
La batterie a une capacité de $500\ \mathrm{Wh}$. Pendant le trajet, le moteur fournit une puissance moyenne de $180\ \mathrm{W}$.

\begin{enumerate}
  \item Rappeler la formule reliant énergie, puissance et durée.
  \item Calculer l'énergie consommée pendant le trajet.
  \item Quelle fraction de la batterie est utilisée ?
  \item En déduire si la batterie suffit pour faire l'aller-retour.
\end{enumerate}
\end{exercice}

\begin{exercice}[PC-LONG-01 -- Partie C : coût de recharge]
Le prix de l'électricité est $0{,}25$ euro par kWh.

\begin{enumerate}
  \item Convertir $500\ \mathrm{Wh}$ en $\mathrm{kWh}$.
  \item Calculer le coût d'une recharge complète.
  \item Comparer ce coût à un ticket de bus à $1{,}80$ euro.
  \item Donner une limite de cette comparaison.
\end{enumerate}
\end{exercice}

\begin{correction}[PC-LONG-01 -- Correction synthétique]
\[
54\ \mathrm{min}=\frac{54}{60}=0{,}9\ \mathrm{h}.
\]
\[
v=\frac{18}{0{,}9}=20\ \mathrm{km/h}.
\]
\[
20\ \mathrm{km/h}=\frac{20}{3{,}6}\approx5{,}6\ \mathrm{m/s}.
\]
Énergie :
\[
E=P\times t=180\times0{,}9=162\ \mathrm{Wh}.
\]
Fraction utilisée :
\[
\frac{162}{500}=0{,}324.
\]
Donc environ $32{,}4\%$ de la batterie.

Une recharge complète :
\[
500\ \mathrm{Wh}=0{,}5\ \mathrm{kWh}.
\]
Coût :
\[
0{,}5\times0{,}25=0{,}125\ \text{euro}.
\]
\end{correction}

\newpage

\subsection{PC-LONG-02 -- Qualité de l'eau, masse volumique et ions}

\begin{objectif}
Durée conseillée : 1 h 10.\\
Chapitres : masse volumique, ions, tests chimiques, concentration simple, analyse critique.
\end{objectif}

\begin{exercice}[PC-LONG-02 -- Partie A : masse volumique]
On prélève un échantillon d'eau salée. Un volume de $250\ \mathrm{mL}$ a une masse de $260\ \mathrm{g}$.

\begin{enumerate}
  \item Convertir $250\ \mathrm{mL}$ en $\mathrm{cm^3}$.
  \item Calculer la masse volumique en $\mathrm{g/cm^3}$.
  \item Comparer avec l'eau pure, de masse volumique environ $1{,}0\ \mathrm{g/cm^3}$.
  \item Proposer une explication.
\end{enumerate}
\end{exercice}

\begin{exercice}[PC-LONG-02 -- Partie B : ions]
On réalise des tests chimiques sur cette eau.
\begin{itemize}
  \item Avec du nitrate d'argent, on observe un précipité blanc.
  \item Avec de la soude, on n'observe pas de précipité coloré.
\end{itemize}

\begin{enumerate}
  \item Quel ion peut être mis en évidence par le nitrate d'argent ?
  \item Que signifie l'apparition d'un précipité ?
  \item Pourquoi le résultat avec la soude ne permet-il pas d'affirmer la présence d'ions cuivre ou fer ?
  \item Donner une conclusion sur la composition probable de l'eau.
\end{enumerate}
\end{exercice}

\begin{exercice}[PC-LONG-02 -- Partie C : évaporation]
On laisse évaporer $100\ \mathrm{mL}$ de cette eau salée. Il reste $3{,}5\ \mathrm{g}$ de solide.

\begin{enumerate}
  \item Quelle masse de solide obtiendrait-on avec $1\ \mathrm{L}$ d'eau salée ?
  \item Exprimer cette concentration massique en $\mathrm{g/L}$.
  \item Peut-on affirmer que le solide est uniquement du sel ? Justifier.
  \item Donner deux sources d'erreur possibles dans l'expérience.
\end{enumerate}
\end{exercice}

\begin{correction}[PC-LONG-02 -- Correction synthétique]
\[
250\ \mathrm{mL}=250\ \mathrm{cm^3}.
\]
Masse volumique :
\[
\rho=\frac{m}{V}=\frac{260}{250}=1{,}04\ \mathrm{g/cm^3}.
\]
L'eau salée est plus dense que l'eau pure.

Le nitrate d'argent met en évidence les ions chlorure $\mathrm{Cl^-}$ par formation d'un précipité blanc.
Avec $100\ \mathrm{mL}$, on obtient $3{,}5$ g de solide. Avec $1\ \mathrm{L}=1000\ \mathrm{mL}$, on obtiendrait :
\[
3{,}5\times10=35\ \mathrm{g}.
\]
La concentration massique est donc :
\[
35\ \mathrm{g/L}.
\]
On ne peut pas affirmer que tout le solide est du sel pur : il peut y avoir d'autres espèces dissoutes.
\end{correction}